% GENERATED FILE. Edit canonical lessons, then run npm run generate:books.
% canonical-source-hash: fnv1a64:a2f80048feb497b6
% canonical-lessons: RU-C235-krichat, RU-C235-sheptat, RU-C235-molchat, RU-C235-ulybatsya, RU-C235-serditsya

\chapter{To shout, To whisper, To be silent, To smile, To be angry}
\label{ch:ru-to-shout-to-whisper-to-be-silent-to-smile-to-be-angry}
\hlchaptermodality{235}

\begin{chapteropening}
\textbf{By the end of this chapter:} \emph{I can say to shout, to whisper, to be silent, to smile and to be angry.}

Complete the last lesson of chapter 235: I can say to shout, to whisper, to be silent, to smile and to be angry.
\end{chapteropening}

\section[\texorpdfstring{\ru{кричать} (krichát)}{кричать (krichát)}]{\ru{кричать} (krichát) — to shout}

\label{lesson:RU-C235-krichat}



\begin{quote}
Before the new one: say the Russian for to grow, then the Russian for to feed.
\end{quote}

\subsection*{You'll want to know: \ru{кричать}}
\textbf{\ru{кричать}} (\emph{krichát}) — \textquotedblleft{}to shout\textquotedblright{}.

\begin{grammarlens}[title={What you've built}]
One more word for this chapter.
\end{grammarlens}

\subsection*{Your turn}
\begin{itemize}
\raggedright
  \item \emph{Say it:} \emph{krichát}
  \item \emph{Say it:} \emph{krichát}, once more
  \item \emph{Say it:} say \emph{krichát}
  \item \emph{Recall:} say the Russian for to grow, then the Russian for to feed, then say \emph{krichát} again
\end{itemize}

\begin{tcolorbox}[breakable,colback=teal!4,colframe=teal!35!black,title={Before you move on}]
What does \ru{кричать} mean? (\textquotedblleft{}To shout\textquotedblright{}.) Say it once more.
\end{tcolorbox}
\section[\texorpdfstring{\ru{шептать} (sheptát)}{шептать (sheptát)}]{\ru{шептать} (sheptát) — to whisper}

\label{lesson:RU-C235-sheptat}



\begin{quote}
Before the new one: say the Russian for to feed, then the Russian for to shout.
\end{quote}

\subsection*{You'll want to know: \ru{шептать}}
\textbf{\ru{шептать}} (\emph{sheptát}) — \textquotedblleft{}to whisper\textquotedblright{}.

\begin{grammarlens}[title={What you've built}]
One more word for this chapter.
\end{grammarlens}

\subsection*{Your turn}
\begin{itemize}
\raggedright
  \item \emph{Say it:} \emph{sheptát}
  \item \emph{Say it:} \emph{sheptát}, once more
  \item \emph{Say it:} say \emph{sheptát}
  \item \emph{Recall:} say the Russian for to feed, then the Russian for to shout, then say \emph{sheptát} again
\end{itemize}

\begin{tcolorbox}[breakable,colback=teal!4,colframe=teal!35!black,title={Before you move on}]
What does \ru{шептать} mean? (\textquotedblleft{}To whisper\textquotedblright{}.) Say it once more.
\end{tcolorbox}
\section[\texorpdfstring{\ru{молчать} (molchát)}{молчать (molchát)}]{\ru{молчать} (molchát) — to be silent}

\label{lesson:RU-C235-molchat}



\begin{quote}
Before the new one: say the Russian for to shout, then the Russian for to whisper.
\end{quote}

\subsection*{You'll want to know: \ru{молчать}}
\textbf{\ru{молчать}} (\emph{molchát}) — \textquotedblleft{}to be silent\textquotedblright{}.

\begin{grammarlens}[title={What you've built}]
One more word for this chapter.
\end{grammarlens}

\subsection*{Your turn}
\begin{itemize}
\raggedright
  \item \emph{Say it:} \emph{molchát}
  \item \emph{Say it:} \emph{molchát}, once more
  \item \emph{Say it:} say \emph{molchát}
  \item \emph{Recall:} say the Russian for to shout, then the Russian for to whisper, then say \emph{molchát} again
\end{itemize}

\begin{tcolorbox}[breakable,colback=teal!4,colframe=teal!35!black,title={Before you move on}]
What does \ru{молчать} mean? (\textquotedblleft{}To be silent\textquotedblright{}.) Say it once more.
\end{tcolorbox}
\section[\texorpdfstring{\ru{улыбаться} (ulybátsya)}{улыбаться (ulybátsya)}]{\ru{улыбаться} (ulybátsya) — to smile}

\label{lesson:RU-C235-ulybatsya}



\begin{quote}
Before the new one: say the Russian for to whisper, then the Russian for to be silent.
\end{quote}

\subsection*{You'll want to know: \ru{улыбаться}}
\textbf{\ru{улыбаться}} (\emph{ulybátsya}) — \textquotedblleft{}to smile\textquotedblright{}.

\begin{grammarlens}[title={What you've built}]
One more word for this chapter.
\end{grammarlens}

\subsection*{Your turn}
\begin{itemize}
\raggedright
  \item \emph{Say it:} \emph{ulybátsya}
  \item \emph{Say it:} \emph{ulybátsya}, once more
  \item \emph{Say it:} say \emph{ulybátsya}
  \item \emph{Recall:} say the Russian for to whisper, then the Russian for to be silent, then say \emph{ulybátsya} again
\end{itemize}

\begin{tcolorbox}[breakable,colback=teal!4,colframe=teal!35!black,title={Before you move on}]
What does \ru{улыбаться} mean? (\textquotedblleft{}To smile\textquotedblright{}.) Say it once more.
\end{tcolorbox}
\section[\texorpdfstring{\ru{сердиться} (serdítsya)}{сердиться (serdítsya)}]{\ru{сердиться} (serdítsya) — to be angry}

\label{lesson:RU-C235-serditsya}



\begin{quote}
Before the new one: say the Russian for to be silent, then the Russian for to smile.
\end{quote}

\subsection*{You'll want to know: \ru{сердиться}}
\textbf{\ru{сердиться}} (\emph{serdítsya}) — \textquotedblleft{}to be angry\textquotedblright{}.

\begin{grammarlens}[title={What you've built}]
One more word for this chapter.
\end{grammarlens}

\subsection*{Your turn}
\begin{itemize}
\raggedright
  \item \emph{Say it:} \emph{serdítsya}
  \item \emph{Say it:} \emph{serdítsya}, once more
  \item \emph{Say it:} say \emph{serdítsya}
  \item \emph{Recall:} say the Russian for to be silent, then the Russian for to smile, then say \emph{serdítsya} again
\end{itemize}

\begin{tcolorbox}[breakable,colback=teal!4,colframe=teal!35!black,title={Before you move on}]
What does \ru{сердиться} mean? (\textquotedblleft{}To be angry\textquotedblright{}.) Say it once more.
\end{tcolorbox}
